\section{Laplace, Stieltjes, and Bernstein transport}
The most portable part of the theory is permanence.  Once a positive kernel, a Stieltjes representation, or a Bernstein derivative has been found, closure under products, mixtures, composition, integration, and stable subordination transports the result to many apparently different source problems.

The section is written as a toolbox with proofs, not as a catalogue.  Definitions and closure principles come before their consequences, bridge propositions explain how a source-specific object enters the common framework, and APP consequences appear only after the transport chain that proves them is visible.
\begin{theorem}
\label{res:cm-closure-product-positive-mixture}
Completely monotone functions on \((0,\infty)\) are closed under pointwise products, nonnegative finite sums, and nonnegative parameter integrals whenever the integral is finite.

\emph{Source context.} This result is used in the cited application chain \cite{mishra-swaminathan-ramanujan}.
\end{theorem}
\begin{proof}
The density
\[
\phi_0(t)=\frac{1}{t(\pi^2+\log^2 t)}
\]
is completely monotone on \((0,\infty)\). Consequently \(I_R\) is a Stieltjes function. Moreover, the antiderivative
\[
\widetilde I_R(x)=a+\int_0^\infty (1-e^{-xt})\phi_0(t)\,dt
\]
is a complete Bernstein function.

For \(u=\log t\),
\[
\int_0^1 e^{-au}\sin(\pi a)\,da
=\frac{\pi(1+e^{-u})}{u^2+\pi^2}.
\]
Substituting \(u=\log t\) gives
\[
\int_0^1 t^{-a}\sin(\pi a)\,da
=\frac{\pi(1+t^{-1})}{\pi^2+\log^2 t}.
\]
Therefore
\[
\phi_0(t)
=\frac{1}{t(\pi^2+\log^2 t)}
=\frac{1}{\pi(1+t)}\int_0^1 t^{-a}\sin(\pi a)\,da.
\]

For \(0<a<1\), \(t^{-a}\) is completely monotone because
\[
t^{-a}=\frac1{\Gamma(a)}\int_0^\infty e^{-ts}s^{a-1}\,ds.
\]
Also \(t\mapsto(1+t)^{-1}\) is completely monotone. Complete monotonicity is closed under products and positive mixtures, and \(\sin(\pi a)>0\) on \((0,1)\). Hence \(\phi_0\) is completely monotone.

If \(\phi\) is completely monotone and
\[
F(x)=\int_0^\infty e^{-xt}\phi(t)\,dt
\]
is finite for \(x>0\), then \(F\) is Stieltjes: write \(\phi(t)=\int_0^\infty e^{-ts}\,d\mu(s)\) and use Tonelli to obtain
\[
F(x)=\int_0^\infty\frac{d\mu(s)}{x+s}.
\]
Applying this to \(\phi_0\) proves that \(I_R\) is Stieltjes.

Finally, the source already proves that
\[
\widetilde I_R(x)=a+\int_0^\infty (1-e^{-xt})\phi_0(t)\,dt
\]
is a Bernstein function. The Levy density \(\phi_0\) is completely monotone, and it satisfies the Bernstein integrability condition
\[
\int_0^\infty (1\wedge t)\phi_0(t)\,dt<\infty.
\]
The standard complete-Bernstein criterion for Levy densities therefore gives that \(\widetilde I_R\) is a complete Bernstein function.
\end{proof}

\begin{proposition}
\label{res:prabhakar-q-density-mixture-formula}
If the transform-normalized Pollard measure \(P^\gamma_{\alpha,\beta}\) has density \(p^\gamma_{\alpha,\beta}\) and \(S_\alpha\) has density \(f_\alpha\), then the non-atomic part of Sibisi's \(Q^\gamma_{\alpha,\beta}(\cdot\mid\lambda)\) has density \(q^\gamma_{\alpha,\beta}(t\mid\lambda)=\int_0^\infty (\lambda r)^{-1/\alpha}f_\alpha(t/(\lambda r)^{1/\alpha})p^\gamma_{\alpha,\beta}(r)\,dr\), with a matching pushed-forward atom at zero if \(P^\gamma_{\alpha,\beta}\) has one.

\emph{Source context.} This result is used in the cited application chain \cite{sibisi-prabhakar}.
\end{proposition}
\begin{proof}
Let \(P\) be a finite positive measure on \([0,\infty)\), let
\[
F(u)=\int_0^\infty e^{-ur}\,dP(r),
\]
let \(0<\alpha<1\), \(\lambda>0\), and let \(S_\alpha\) be positive \(\alpha\)-stable:
\[
\mathbb E e^{-xS_\alpha}=e^{-x^\alpha}.
\]
Define a finite positive measure \(Q\) by
\[
Q(A)=\int_{[0,\infty)}
\Pr\{(\lambda r)^{1/\alpha}S_\alpha\in A\}\,dP(r).
\]
Then Tonelli's theorem gives, for \(x>0\),
\[
\int_0^\infty e^{-xt}\,dQ(t)
=\int_{[0,\infty)}
\mathbb E e^{-x(\lambda r)^{1/\alpha}S_\alpha}\,dP(r)
=\int_{[0,\infty)}e^{-\lambda r x^\alpha}\,dP(r)
=F(\lambda x^\alpha).
\]

Apply the lemma to Sibisi's transform-normalized \(P=P^\gamma_{\alpha,\beta}\) and
\[
F(u)=E^\gamma_{\alpha,\beta}(-u).
\]
For \(0<\alpha<1\), \(\gamma>0\), \(\beta>\alpha\gamma\), and \(\lambda>0\), set
\[
Q^\gamma_{\alpha,\beta}(A\mid\lambda)
=\int_{[0,\infty)}
\Pr\{(\lambda r)^{1/\alpha}S_\alpha\in A\}\,
dP^\gamma_{\alpha,\beta}(r).
\]
Then
\[
\int_0^\infty e^{-xt}\,dQ^\gamma_{\alpha,\beta}(t\mid\lambda)
=E^\gamma_{\alpha,\beta}(-\lambda x^\alpha).
\]
By uniqueness of finite Laplace transforms, this is Sibisi's \(Q^\gamma_{\alpha,\beta}(\cdot\mid\lambda)\). \cite{sibisi-prabhakar}.

If \(S_\alpha\) has density \(f_\alpha\), then the non-atomic density part is
\[
q^\gamma_{\alpha,\beta}(t\mid\lambda)
=\int_{(0,\infty)}
(\lambda r)^{-1/\alpha}
f_\alpha\!\left(\frac{t}{(\lambda r)^{1/\alpha}}\right)
dP^\gamma_{\alpha,\beta}(r),
\]
with an atom \(P^\gamma_{\alpha,\beta}(\{0\})\delta_0\) if \(P^\gamma_{\alpha,\beta}\) has an atom at zero. When \(P^\gamma_{\alpha,\beta}\) is represented by Sibisi's Pollard density \(p^\gamma_{\alpha,\beta}(r)\), this becomes
\[
q^\gamma_{\alpha,\beta}(t\mid\lambda)
=\int_0^\infty
(\lambda r)^{-1/\alpha}
f_\alpha\!\left(\frac{t}{(\lambda r)^{1/\alpha}}\right)
p^\gamma_{\alpha,\beta}(r)\,dr.
\]

Because
\[
E^\gamma_{\alpha,\beta}(0)=\frac1{\Gamma(\beta)},
\]
the representing measures \(P^\gamma_{\alpha,\beta}\) and \(Q^\gamma_{\alpha,\beta}(\cdot\mid\lambda)\) have total mass \(1/\Gamma(\beta)\). They are finite positive measures, not generally probability measures. The normalized probability version is
\[
\widehat Q^\gamma_{\alpha,\beta}=\Gamma(\beta)Q^\gamma_{\alpha,\beta},
\qquad
\mathcal L\widehat Q^\gamma_{\alpha,\beta}(x)
=\Gamma(\beta)E^\gamma_{\alpha,\beta}(-\lambda x^\alpha).
\]

The boundary \(\beta=\alpha\gamma\) is not included here because Sibisi's ordinary convolution density uses \(1/\Gamma(\beta-\alpha\gamma)\). A limiting treatment is a separate open frontier.
\end{proof}

\begin{proposition}
\label{res:ks-lemma1-stieltjes-normal-form}
Under the Karp--Sitnik positive-density hypotheses \(b_i>a_i>0\), the function \(F(x)={}_{q+1}F_q(1,a_1,\ldots,a_q;b_1,\ldots,b_q;-x)\) is a normalized Stieltjes function on \((0,\infty)\).
\end{proposition}
\begin{proof}
Let \(q\ge1\) and \(b_i>a_i>0\). Define
\[
F(x)={}_{q+1}F_q(1,a_1,\ldots,a_q;b_1,\ldots,b_q;-x)
\]
and
\[
R(x)=
\frac{{}_{q+1}F_q(1,a_1+1,\ldots,a_q+1;b_1+1,\ldots,b_q+1;-x)}
{{}_{q+1}F_q(1,a_1,\ldots,a_q;b_1,\ldots,b_q;-x)}.
\]
Then \(R\) is a Stieltjes function on \((0,\infty)\), hence completely monotone. Moreover
\[
xR(x)
\]
is a complete Bernstein function up to the positive scalar \(\prod_i b_i/a_i\).

For \(q=1\), this gives the Gauss/Beta slice:
\[
\frac{{}_2F_1(1,a+1;b+1;-x)}{{}_2F_1(1,a;b;-x)}
\]
is Stieltjes and completely monotone for \(b>a>0\).

Karp--Sitnik's positive-density representation gives a compact-support Stieltjes form for \(F\):
\[
F(x)=\int_0^1 \frac{d\mu(s)}{1+sx},
\qquad \mu\ge0,
\qquad \mu([0,1])=F(0)=1.
\]
Thus \(F\) is a nonzero Stieltjes function.

By the standard Stieltjes/complete-Bernstein duality, if \(F\) is nonzero Stieltjes, then
\[
G(x)=\frac1{F(x)}
\]
is a complete Bernstein function. Since \(G(0)=1\), the complete Bernstein representation gives
\[
G(x)=1+\alpha x+\int_0^\infty \frac{x}{x+t}\,d\nu(t),
\qquad \alpha\ge0,\quad \nu\ge0.
\]
Therefore
\[
\frac{G(x)-1}{x}
=\alpha+\int_0^\infty \frac{1}{x+t}\,d\nu(t)
\]
is Stieltjes.

Karp--Sitnik's contiguous defect identity for \(\sigma=1,\delta=1\) is
\[
1-F(x)
=x\,{}_{q+1}F_q(1,a_1+1,\ldots,a_q+1;b_1+1,\ldots,b_q+1;-x)
\prod_{i=1}^q\frac{a_i}{b_i}.
\]
Hence
\[
R(x)=
\prod_{i=1}^q\frac{b_i}{a_i}\,
\frac{1-F(x)}{xF(x)}
=
\prod_{i=1}^q\frac{b_i}{a_i}\,
\frac{G(x)-1}{x}.
\]
A positive scalar multiple of a Stieltjes function is Stieltjes. Thus \(R\) is Stieltjes and completely monotone.

Finally,
\[
xR(x)=\prod_{i=1}^q\frac{b_i}{a_i}\,(G(x)-1).
\]
Since subtracting the finite value \(G(0)=1\) from a complete Bernstein function leaves the same nonnegative drift and Levy measure part, \(G(x)-1\) is complete Bernstein. Hence \(xR(x)\) is complete Bernstein up to the stated positive scalar.

This is a strong local theorem and a reusable Stieltjes/CBF bridge. The subsequent primary-source wording audit did not find an explicit Karp--Sitnik open problem asking for the CM/Stieltjes/CBF upgrade, so this result is bridge-only unless a separate source asks for exactly this strengthening.
\end{proof}

\begin{proposition}
\label{res:jovanovic-treml-arctan-logderivative-cm}
The function \(x\mapsto ((1+x^2)\arctan x)^{-1}\) is completely monotone on \((0,\infty)\), by the Jovanovic--Treml positive-kernel theorem.
\end{proposition}
\begin{proof}
If \(1/\arctan x\) were a Stieltjes function, then its reciprocal \(\arctan x\) would be a Bernstein function. A Bernstein function has completely monotone derivative.

But
\[
(\arctan x)'=\frac1{1+x^2},
\]
and
\[
\left(\frac1{1+x^2}\right)''
=\frac{2(3x^2-1)}{(1+x^2)^3},
\]
which is negative for \(0<x<1/\sqrt3\). Hence \((\arctan x)'\) is not completely monotone, \(\arctan x\) is not Bernstein, and \(1/\arctan x\) is not Stieltjes.
\end{proof}

\begin{proposition}[Heat-flow \(L^2\) energy]
\label{res:heat-flow-l2-energy-completely-monotone}
For every probability measure \(\mu\) on \(\mathbb R^d\), the heat-flow \(L^2\) energy \(N_2(t)=\int_{\mathbb R^d}(G_t*\mu)(x)^2\,dx\) is completely monotone on \((0,\infty)\).
\end{proposition}
\begin{proof}
Let \(\mu\) be a probability measure on \(\mathbb R^d\), and let \(G_t\) be the heat kernel with Fourier transform
\[
\widehat G_t(\xi)=e^{-t|\xi|^2}.
\]
For \(p_t=G_t*\mu\), define
\[
N_2(t)=\int_{\mathbb R^d}p_t(x)^2\,dx,\qquad t>0.
\]
Then \(N_2\) is completely monotone on \((0,\infty)\). If \(S_2(t)=1-N_2(t)\), then \(S_2'\) is completely monotone on \((0,\infty)\). Hence, for each \(t_0>0\), the normalized increment \(S_2(t)-S_2(t_0)\) is Bernstein as a function of \(t\ge t_0\).

Since \(|\widehat\mu(\xi)|\le1\), the function \(e^{-t|\xi|^2}\widehat\mu(\xi)\) is in \(L^2(\mathbb R^d)\) for every \(t>0\). Thus \(p_t\in L^2\), and Plancherel gives
\[
N_2(t)
=(2\pi)^{-d}\int_{\mathbb R^d}e^{-2t|\xi|^2}|\widehat\mu(\xi)|^2\,d\xi.
\]
Define \(\nu_\mu\) to be the locally finite positive pushforward of
\[
(2\pi)^{-d}|\widehat\mu(\xi)|^2\,d\xi
\]
under the map \(\xi\mapsto|\xi|^2\). Then
\[
N_2(t)=\int_0^\infty e^{-2tr}\,d\nu_\mu(r).
\]
The integral is finite for every \(t>0\), because \(|\widehat\mu|\le1\) and \(\int_{\mathbb R^d}e^{-2t|\xi|^2}\,d\xi<\infty\).

For \(k\ge0\), differentiation under the integral is justified at \(t>0\) by the Gaussian factor. It yields
\[
(-1)^kN_2^{(k)}(t)
=\int_0^\infty (2r)^k e^{-2tr}\,d\nu_\mu(r)\ge0.
\]
Therefore \(N_2\) is completely monotone.

Now \(S_2'(t)=-N_2'(t)\), and
\[
S_2'(t)=\int_0^\infty 2r e^{-2tr}\,d\nu_\mu(r),
\]
another positive Laplace transform. Hence \(S_2'\) is completely monotone. This implies that every normalized increment \(S_2(t)-S_2(t_0)\), \(t\ge t_0>0\), is a Bernstein function on that restricted half-line.
\end{proof}

\begin{proposition}
\label{res:bf-reciprocal-is-cm}
If \(g\) is a nonzero Bernstein function on \((0,\infty)\), then \(1/g\) is completely monotone.
\end{proposition}
\begin{proof}
For
\[
f_{\alpha,\beta}(x)=x^{-\alpha}(1+x^\beta)^{-1},\qquad x>0,
\]
the following hold.

1. If \(\alpha\ge0\) and \(0\le\beta\le1\), then \(f_{\alpha,\beta}\) is completely monotone.
2. If \(1<\beta\le2\) and \(\alpha\ge\beta/2\), then \(f_{\alpha,\beta}\) is completely monotone.
3. If \(\beta=2\), then \(f_{\alpha,2}\) is completely monotone if and only if \(\alpha\ge1\).

Consequently, for the Dagum threshold
\[
c(\beta)=\inf\{\alpha\ge0:f_{\alpha,\beta}\in CM\},
\]
one has \(c(1)=0\), \(c(2)=1\), and \(c(\beta)\le\beta/2\) for \(1<\beta\le2\).

For \(0<\beta\le1\), \(x^\beta\) is a Bernstein function, so \(1+x^\beta\) is Bernstein. A nonzero Bernstein function has completely monotone reciprocal. Hence \((1+x^\beta)^{-1}\) is completely monotone. The factor \(x^{-\alpha}\) is completely monotone for \(\alpha\ge0\), and products of completely monotone functions are completely monotone. The case \(\beta=0\) is immediate because \(f_{\alpha,0}(x)=2^{-1}x^{-\alpha}\).

For \(1<\beta\le2\), put \(\gamma=\beta/2\in(1/2,1]\). The elementary Laplace identity
\[
\frac1{x(1+x^2)}=\int_0^\infty e^{-xt}(1-\cos t)\,dt
\]
shows that \(h(x)=1/(x(1+x^2))\) is completely monotone. Since \(x^\gamma\) is Bernstein, \(h(x^\gamma)\) is completely monotone. But
\[
h(x^\gamma)=\frac1{x^\gamma(1+x^{2\gamma})}
=\frac1{x^{\beta/2}(1+x^\beta)}
=f_{\beta/2,\beta}(x).
\]
Multiplication by \(x^{-(\alpha-\beta/2)}\) proves complete monotonicity for \(\alpha\ge\beta/2\).

For \(\beta=2\), sufficiency for \(\alpha\ge1\) follows from the preceding paragraph. For necessity, first \(f_{0,2}(x)=1/(1+x^2)\) is not completely monotone, for example because its inverse Laplace kernel is \(\sin t\).

Let \(0<\alpha<1\). The inverse Laplace kernel of \(x^{-\alpha}\) is \(t^{\alpha-1}/\Gamma(\alpha)\), and the inverse Laplace kernel of \(1/(1+x^2)\) is \(\sin t\). Hence the inverse kernel of \(f_{\alpha,2}\) is
\[
\eta_\alpha(t)=\frac1{\Gamma(\alpha)}\int_0^t(t-s)^{\alpha-1}\sin s\,ds.
\]
At \(t=2\pi\),
\[
\Gamma(\alpha)\eta_\alpha(2\pi)
=\int_0^\pi\left((2\pi-s)^{\alpha-1}-(\pi-s)^{\alpha-1}\right)\sin s\,ds.
\]
Since \(\alpha-1<0\), the bracket is negative for \(0<s<\pi\). Therefore \(\eta_\alpha(2\pi)<0\). By uniqueness of the Laplace transform for locally finite measures, \(f_{\alpha,2}\) cannot be completely monotone. Thus \(f_{\alpha,2}\in CM\) exactly when \(\alpha\ge1\).
\end{proof}

\begin{proposition}
\label{res:baskakov-alpha1-r6-positive-density-factorization}
For \(m=3\), equivalently \(r=6\), the Baskakov \(\alpha=1\) function \(f^{[6]}_1(x)=1/((1+x)^6-x^6)\) has inverse-Laplace density \(\rho_3(t)=\frac{4}{3}e^{-t/2}(1-\cos(t/(2\sqrt3)))^2(2+\cos(t/(2\sqrt3)))\), hence is completely monotone.
\end{proposition}
\begin{proof}
Set
\[
D_m(x)=(1+x)^{2m}-x^{2m}.
\]
The roots are obtained from
\[
\frac{1+x}{x}=\omega_k,\qquad
\omega_k=e^{\pi i k/m},\qquad k=1,\ldots,2m-1.
\]
Thus
\[
\lambda_{k,m}
=\frac{1}{\omega_k-1}
=-\frac12-\frac{i}{2}\cot\frac{\pi k}{2m}.
\]
At this root,
\[
D_m'(\lambda_{k,m})
=2m\lambda_{k,m}^{2m-1}(\omega_k^{-1}-1),
\]
and direct simplification gives the real residue
\[
R_{k,m}
=\frac{1}{D_m'(\lambda_{k,m})}
=\frac{(-1)^{m+k}}{2m}
\left(2\sin\frac{\pi k}{2m}\right)^{2m-2}.
\]
Consequently the inverse-Laplace density is
\[
\rho_m(t)
=\sum_{k=1}^{2m-1}R_{k,m}e^{\lambda_{k,m}t}.
\]
Pairing conjugate roots gives
\[
\rho_m(t)
=\frac{e^{-t/2}}{2m}
\left[
2^{2m-2}
+2\sum_{k=1}^{m-1}
(-1)^{m+k}
\left(2\sin\frac{\pi k}{2m}\right)^{2m-2}
\cos\!\left(
\frac{t}{2}\cot\frac{\pi k}{2m}
\right)
\right].
\]

This proves the finite trigonometric density formula used by the diagnostic
candidate.

For \(m=2\), this gives
\[
\rho_2(t)=e^{-t/2}\left(1-\cos\frac t2\right)
=2e^{-t/2}\sin^2\frac t4\ge0,
\]
recovering the previously admitted \(r=4,\alpha=1\) seed.

For \(m=3\),
\[
\rho_3(t)
=\frac{e^{-t/2}}6
\left[
16+2\cos\left(\frac{\sqrt3\,t}{2}\right)
-18\cos\left(\frac{t}{2\sqrt3}\right)
\right].
\]
Set \(u=t/(2\sqrt3)\).  Since
\(\cos(3u)=4\cos^3u-3\cos u\),
\[
\rho_3(t)
=\frac{4}{3}e^{-t/2}(1-\cos u)^2(2+\cos u)\ge0.
\]
Thus \(r=6,\alpha=1\) is completely monotone by an explicit positive density.

For \(m=4\), write \(h_4(t)=e^{t/2}\rho_4(t)\).  The density formula gives
\[
h_4(t)
=8+2\cos\frac t2
-\frac{(2-\sqrt2)^3}{4}
\cos\left(\frac{1+\sqrt2}{2}t\right)
-\frac{(2+\sqrt2)^3}{4}
\cos\left(\frac{\sqrt2-1}{2}t\right).
\]
At \(t=10\pi\),
\[
\cos(t/2)=\cos(5\pi)=-1,
\]
and both other cosine terms equal \(-\cos(5\pi\sqrt2)\).  Since
\[
(2-\sqrt2)^3+(2+\sqrt2)^3=40,
\]
we obtain
\[
h_4(10\pi)=6+10\cos(5\pi\sqrt2).
\]
Now
\[
0<5\sqrt2-7<\frac14
\]
because \(7/5<\sqrt2<29/20\).  Hence, with
\(\delta=5\sqrt2-7\),
\[
\cos(5\pi\sqrt2)=\cos(7\pi+\pi\delta)=-\cos(\pi\delta)
<-\cos\frac\pi4=-\frac{\sqrt2}{2}<-\frac35.
\]
Therefore
\[
h_4(10\pi)<6+10\left(-\frac35\right)=0,
\]
and
\[
\rho_4(10\pi)=e^{-5\pi}h_4(10\pi)<0.
\]
Since the inverse-Laplace transform is unique for measures on
\([0,\infty)\), \(f^{[8]}_1(x)=1/((1+x)^8-x^8)\) cannot be completely monotone.
\end{proof}

\begin{proposition}
\label{res:bazhlekova-odd-derivative-polynomial-normal-form}
For \(h(s)=\sqrt{c s^a+d s^b}\), with \(c,d>0\), \(\Delta=a-b\), \(B=b/2\), and \(y=(c/d)s^\Delta\), the derivatives have the form \(h^{(n)}(s)=\sqrt d\,s^{B-n}(1+y)^{1/2-n}Q_n(y)\), where \(Q_0=1\) and \(Q_{n+1}=(B-n)(1+y)Q_n+\Delta y((1+y)Q_n'+(1/2-n)Q_n)\).
\end{proposition}
\begin{proof}
If for an odd \(n=2q+1\) there is \(y_0>0\) such that
\[
Q_n(y_0)<0,
\]
then for every \(c,d>0\) the point
\[
s_0=\left(\frac{d y_0}{c}\right)^{1/(a-b)}
\]
satisfies \(h^{(n)}(s_0)<0\).

A Bernstein function \(h\) has completely monotone derivative, hence for odd \(n\ge3\) it must satisfy \(h^{(n)}\ge0\). Thus the displayed sign proves that \(h\) is not Bernstein.

More directly for the Bazhlekova transform,
\[
F_x(s)=e^{-x h(s)}=\mathcal L\{w_t(x,\cdot)\}(s),
\]
we have
\[
F_x^{(n)}(s_0)=-x h^{(n)}(s_0)+O(x^2).
\]
For odd \(n\), complete monotonicity would require \(F_x^{(n)}(s_0)\le0\). If \(h^{(n)}(s_0)<0\), then \(F_x^{(n)}(s_0)>0\) for all sufficiently small \(x>0\), so \(w_t(x,\cdot)\) cannot be nonnegative.

The previous exact inner-gap example is recovered by the recurrence. For
\[
(a,b)=\left(\frac{28}{25},\frac1{50}\right),
\]
one has
\[
a-b=\frac{11}{10},\qquad
(a-1)^2+(b-1)^2=\frac{2437}{2500}<1,
\]
and
\[
Q_5(1)=-\frac{5570045943}{10000000000}<0.
\]
This matches the earlier value
\[
h^{(5)}(1)
=-\frac{5570045943\sqrt2}{320000000000}<0
\]
when \(c=d=1\).

The seventh derivative adds a new exact residual seed. For
\[
(a,b)=\left(\frac{107}{100},\frac1{100}\right),
\]
the point lies in the residual inner gap and exact Sturm counting gives no positive-root obstruction for \(Q_5\). But
\[
Q_7\left(\frac32\right)
=-\frac{56394696198326721417}{327680000000000000}<0.
\]

The ninth derivative adds another seed. For
\[
(a,b)=\left(\frac{53}{50},\frac1{100}\right),
\]
exact Sturm counting gives no positive-root obstruction for \(Q_5\) or \(Q_7\), but
\[
Q_9\left(\frac74\right)
=-\frac{3262653919060921278409379787}
{262144000000000000000000}<0.
\]

The finite \(5,7,9\) diagnostic does not cover the residual inner gap. At
\[
(a,b)=\left(\frac32,\frac25\right)
\]
we have
\[
0<b<a-1,\qquad a-b=\frac{11}{10}>1,\qquad
(a-1)^2+(b-1)^2=\frac{61}{100}<1.
\]
Exact Sturm counts show that \(Q_5,Q_7,Q_9\) have no positive roots. Since each has positive value at \(0\), each is positive on \(y>0\). Thus the fifth, seventh, and ninth derivative tests cannot decide this exponent pair.

The same noncoverage phenomenon also occurs at the near-boundary residual point
\[
(a,b)=\left(\frac{11}{10},\frac1{20}\right),
\qquad
(a-1)^2+(b-1)^2=\frac{73}{80}<1.
\]
Again \(Q_5,Q_7,Q_9\) have no positive roots and are positive at \(0\).

The derivative-polynomial normal form for \(h^{(n)}\).
The odd-derivative small-\(x\) obstruction criterion.
The exact fifth/seventh/ninth negative seeds above.
The exact finite \(5,7,9\) no-cover seeds above.
\end{proof}

The next result applies the preceding method to the fixed source problem \textbf{APP-0014} \cite{mishra-swaminathan-ramanujan}.
\begin{corollary}[APP-0014: Ramanujan integral is Stieltjes and has a complete Bernstein primitive]
\label{res:ramanujan-density-logsquare-cm}
The function \(\phi_0(t)=1/[t(\pi^2+\log^2 t)]\) is completely monotone on \((0,\infty)\).

This resolves the fixed application label \textbf{APP-0014}: Classify the Ramanujan integral by a standard positive-kernel transform rather than only by inequalities.
\emph{Source reference.} \cite{mishra-swaminathan-ramanujan}.

\emph{Dependency.} The proof uses the preceding result(s) \ref{res:cm-closure-product-positive-mixture}.
\end{corollary}
\begin{proof}
For \(t>0\),
\[
 {1\over t(\pi^2+\log^2t)}
   ={1\over \pi(1+t)}\int_0^1 t^{-u}\sin(\pi u)\,du .
\]
The right side is a positive superposition of completely monotone powers and is
therefore completely monotone as a function of \(t\).  The Laplace transform of
a completely monotone density is Stieltjes.  Integrating the same density
against \(1-e^{-xt}\) gives a complete Bernstein function by the standard
Stieltjes-density characterization of complete Bernstein functions.
\end{proof}

\begin{proposition}
\label{res:ks-general-q-delta1-sigma1-ratio-stieltjes-cm}
For \(q\ge1\) and \(b_i>a_i>0\), the shifted ratio \({}_{q+1}F_q(1,a_1+1,\ldots,a_q+1;b_1+1,\ldots,b_q+1;-x)/{}_{q+1}F_q(1,a_1,\ldots,a_q;b_1,\ldots,b_q;-x)\) is a Stieltjes function, hence completely monotone, on \((0,\infty)\).

\emph{Dependency.} The proof uses the preceding result(s) \ref{res:cbf-difference-quotient-stieltjes}.
\end{proposition}
\begin{proof}
Let \(q\ge1\) and \(b_i>a_i>0\). Define
\[
F(x)={}_{q+1}F_q(1,a_1,\ldots,a_q;b_1,\ldots,b_q;-x)
\]
and
\[
R(x)=
\frac{{}_{q+1}F_q(1,a_1+1,\ldots,a_q+1;b_1+1,\ldots,b_q+1;-x)}
{{}_{q+1}F_q(1,a_1,\ldots,a_q;b_1,\ldots,b_q;-x)}.
\]
Then \(R\) is a Stieltjes function on \((0,\infty)\), hence completely monotone. Moreover
\[
xR(x)
\]
is a complete Bernstein function up to the positive scalar \(\prod_i b_i/a_i\).

For \(q=1\), this gives the Gauss/Beta slice:
\[
\frac{{}_2F_1(1,a+1;b+1;-x)}{{}_2F_1(1,a;b;-x)}
\]
is Stieltjes and completely monotone for \(b>a>0\).

Karp--Sitnik's positive-density representation gives a compact-support Stieltjes form for \(F\):
\[
F(x)=\int_0^1 \frac{d\mu(s)}{1+sx},
\qquad \mu\ge0,
\qquad \mu([0,1])=F(0)=1.
\]
Thus \(F\) is a nonzero Stieltjes function.

By the standard Stieltjes/complete-Bernstein duality, if \(F\) is nonzero Stieltjes, then
\[
G(x)=\frac1{F(x)}
\]
is a complete Bernstein function. Since \(G(0)=1\), the complete Bernstein representation gives
\[
G(x)=1+\alpha x+\int_0^\infty \frac{x}{x+t}\,d\nu(t),
\qquad \alpha\ge0,\quad \nu\ge0.
\]
Therefore
\[
\frac{G(x)-1}{x}
=\alpha+\int_0^\infty \frac{1}{x+t}\,d\nu(t)
\]
is Stieltjes.

Karp--Sitnik's contiguous defect identity for \(\sigma=1,\delta=1\) is
\[
1-F(x)
=x\,{}_{q+1}F_q(1,a_1+1,\ldots,a_q+1;b_1+1,\ldots,b_q+1;-x)
\prod_{i=1}^q\frac{a_i}{b_i}.
\]
Hence
\[
R(x)=
\prod_{i=1}^q\frac{b_i}{a_i}\,
\frac{1-F(x)}{xF(x)}
=
\prod_{i=1}^q\frac{b_i}{a_i}\,
\frac{G(x)-1}{x}.
\]
A positive scalar multiple of a Stieltjes function is Stieltjes. Thus \(R\) is Stieltjes and completely monotone.

Finally,
\[
xR(x)=\prod_{i=1}^q\frac{b_i}{a_i}\,(G(x)-1).
\]
Since subtracting the finite value \(G(0)=1\) from a complete Bernstein function leaves the same nonnegative drift and Levy measure part, \(G(x)-1\) is complete Bernstein. Hence \(xR(x)\) is complete Bernstein up to the stated positive scalar.

This is a strong local theorem and a reusable Stieltjes/CBF bridge. The subsequent primary-source wording audit did not find an explicit Karp--Sitnik open problem asking for the CM/Stieltjes/CBF upgrade, so this result is bridge-only unless a separate source asks for exactly this strengthening.
\end{proof}

\begin{proposition}[reciprocal arctan is logarithmically completely monotone]
\label{res:reciprocal-arctan-lcm}
The function \(x\mapsto1/\arctan x\) is logarithmically completely monotone on \((0,\infty)\).

\emph{Dependency.} The proof uses the preceding result(s) \ref{res:jovanovic-treml-arctan-logderivative-cm}.
\end{proposition}
\begin{proof}
If \(1/\arctan x\) were a Stieltjes function, then its reciprocal \(\arctan x\) would be a Bernstein function. A Bernstein function has completely monotone derivative.

But
\[
(\arctan x)'=\frac1{1+x^2},
\]
and
\[
\left(\frac1{1+x^2}\right)''
=\frac{2(3x^2-1)}{(1+x^2)^3},
\]
which is negative for \(0<x<1/\sqrt3\). Hence \((\arctan x)'\) is not completely monotone, \(\arctan x\) is not Bernstein, and \(1/\arctan x\) is not Stieltjes.
\end{proof}

\begin{proposition}
\label{res:bs-c-half-canonical-r-sequence-cm}
For the Bondesson--Steutel distribution at \(c=1/2\), the canonical sequence \(r_n=(n+1/2)P_n(1/2)\) is completely monotone, with \(r_n=\frac{1}{2\pi}\int_0^1 x^{n+1/2}(1-x)^{-1/2}\,dx\).

\emph{Dependency.} The proof uses the preceding result(s) \ref{res:bs-c-half-catalan-hausdorff-representation}.
\end{proposition}
\begin{proof}
For \(c=1/2\), the source pgf is
\[
P(z)=\frac{1-\sqrt{1-z}}{z}.
\]
Using the Catalan generating function
\[
\sum_{n\ge0}C_n x^n=\frac{1-\sqrt{1-4x}}{2x},
\]
with \(x=z/4\), we get
\[
P(z)=\frac12\sum_{n\ge0}C_n\left(\frac z4\right)^n.
\]
Hence
\[
P_n(1/2)=\frac{C_n}{2\cdot4^n}
=\frac{1}{2\cdot4^n(n+1)}\binom{2n}{n}.
\]

Now
\[
\frac1\pi\int_0^1 x^{n-1/2}(1-x)^{1/2}\,dx
=\frac1\pi B(n+1/2,3/2).
\]
Since \(\Gamma(3/2)=\sqrt\pi/2\) and
\[
\Gamma(n+1/2)=\frac{(2n)!\sqrt\pi}{4^n n!},
\]
this beta integral equals
\[
\frac{(2n)!}{2\cdot4^n n!(n+1)!}
=\frac{C_n}{2\cdot4^n}
=P_n(1/2).
\]
Therefore
\[
P_n(1/2)=\frac1\pi\int_0^1 x^{n-1/2}(1-x)^{1/2}\,dx.
\]
This is a Hausdorff moment representation on \([0,1]\), so \((P_n(1/2))\) is completely monotone as a sequence.

For \(c=1/2\), the canonical sequence is
\[
r_n=(n+1/2)P_n(1/2).
\]
Using the Hausdorff representation,
\[
r_n=\frac{n+1/2}{\pi}\int_0^1 x^{n-1/2}(1-x)^{1/2}\,dx.
\]
Because
\[
\frac{d}{dx}x^{n+1/2}=(n+1/2)x^{n-1/2},
\]
integration by parts gives
\[
r_n
=\frac1\pi\int_0^1 (1-x)^{1/2}\,d(x^{n+1/2})
=\frac{1}{2\pi}\int_0^1 x^{n+1/2}(1-x)^{-1/2}\,dx.
\]
The boundary terms vanish at \(0\) and \(1\). Thus
\[
r_n=\int_0^1 x^n\,d\nu(x),
\qquad
d\nu(x)=\frac{1}{2\pi}x^{1/2}(1-x)^{-1/2}\,dx.
\]
This is a positive finite measure on \([0,1]\). Hence \((r_n)\) is completely monotone.
\end{proof}

\begin{proposition}[Ramanujan integral I\_R is Stieltjes]
\label{res:ramanujan-integral-stieltjes}
The Ramanujan integral \(I_R(x)=\int_0^\infty e^{-xt}\,dt/[t(\pi^2+\log^2 t)]\) is a Stieltjes function on \((0,\infty)\).

\emph{Dependency.} The proof uses the preceding result(s) \ref{res:ramanujan-density-logsquare-cm}.
\end{proposition}
\begin{proof}
The density
\[
\phi_0(t)=\frac{1}{t(\pi^2+\log^2 t)}
\]
is completely monotone on \((0,\infty)\). Consequently \(I_R\) is a Stieltjes function. Moreover, the antiderivative
\[
\widetilde I_R(x)=a+\int_0^\infty (1-e^{-xt})\phi_0(t)\,dt
\]
is a complete Bernstein function.

For \(u=\log t\),
\[
\int_0^1 e^{-au}\sin(\pi a)\,da
=\frac{\pi(1+e^{-u})}{u^2+\pi^2}.
\]
Substituting \(u=\log t\) gives
\[
\int_0^1 t^{-a}\sin(\pi a)\,da
=\frac{\pi(1+t^{-1})}{\pi^2+\log^2 t}.
\]
Therefore
\[
\phi_0(t)
=\frac{1}{t(\pi^2+\log^2 t)}
=\frac{1}{\pi(1+t)}\int_0^1 t^{-a}\sin(\pi a)\,da.
\]

For \(0<a<1\), \(t^{-a}\) is completely monotone because
\[
t^{-a}=\frac1{\Gamma(a)}\int_0^\infty e^{-ts}s^{a-1}\,ds.
\]
Also \(t\mapsto(1+t)^{-1}\) is completely monotone. Complete monotonicity is closed under products and positive mixtures, and \(\sin(\pi a)>0\) on \((0,1)\). Hence \(\phi_0\) is completely monotone.

If \(\phi\) is completely monotone and
\[
F(x)=\int_0^\infty e^{-xt}\phi(t)\,dt
\]
is finite for \(x>0\), then \(F\) is Stieltjes: write \(\phi(t)=\int_0^\infty e^{-ts}\,d\mu(s)\) and use Tonelli to obtain
\[
F(x)=\int_0^\infty\frac{d\mu(s)}{x+s}.
\]
Applying this to \(\phi_0\) proves that \(I_R\) is Stieltjes.

Finally, the source already proves that
\[
\widetilde I_R(x)=a+\int_0^\infty (1-e^{-xt})\phi_0(t)\,dt
\]
is a Bernstein function. The Levy density \(\phi_0\) is completely monotone, and it satisfies the Bernstein integrability condition
\[
\int_0^\infty (1\wedge t)\phi_0(t)\,dt<\infty.
\]
The standard complete-Bernstein criterion for Levy densities therefore gives that \(\widetilde I_R\) is a complete Bernstein function.
\end{proof}
